% s03-1-2-characters.tex —— 3.1.2 小节：AB 的主角（印刷页 61–69）

\subsection{AB 的主角}
\label{sec:ab-characters}

本小节及随后两小节回顾或定义全文将反复出现的 $\FactAlgCat$ 与 $\FactModCat_{x_0}$ 中的
各种对象。

先回顾 \cite[1.3 节]{ABC+24a} 中给出的因子化范畴 $\Whit^{!}(G)$ 的定义。

为此先考虑 \emph{Ran Grassmann 叠} $\Gr_{G,\Ran}$。对试验仿射概形 $S$ 及
$x\colon S\to \Ran$，$x$ 到 $\Gr_{G,\Ran}$ 的提升由一对
\[
(P_G,\alpha)
\]
给出：$S\times X$ 上的主 $G$-丛 $P_G$，以及在 $S\times X$ 中 $x$ 的图像之补集上
$P_G$ 的平凡化
\[
\alpha\colon P_G|_{(S\times X)\setminus \Gamma_x}\xrightarrow{\ \sim\ }P_G^{\mathrm{triv}}.
\]
\srcpage{61}

Ran Grassmann 叠自然带有（带单位的）因子化空间结构，于是得到（带单位的）因子化范畴
$D\text{-}\mathrm{mod}_{1/2}(\Gr_G)$（临界扭转 $D$-模的讨论见 \cite[1.1 节]{ABC+24a}），
其定义为
\[
D\text{-}\mathrm{mod}_{1/2}(\Gr_G)_S
=D\text{-}\mathrm{mod}_{1/2}(S\times_{\Ran}\Gr_{G,\Ran}).
\]
此外，因子化群 ind-概形 $LG_{\Ran}$（从而因子化子群 ind-概形 $LN_{\Ran}$）作用在
$\Gr_{G,\Ran}$ 上。现在基本上要做的是把 $\Whit^{!}(G)$ 定义为
$D\text{-}\mathrm{mod}_{1/2}(\Gr_G)$ 的 Whittaker 不变量范畴，即
\[
D\text{-}\mathrm{mod}_{1/2}(\Gr_G)^{(LN,\chi)},
\]
更精确地说，
\[
S\mapsto D\text{-}\mathrm{mod}_{1/2}(\Gr_G)^{(LN^S,\chi^S)}_S .
\]
但还必须纳入 \cite[1.2 节]{ABC+24a} 中的几何扭转构造，这里简略复述。

设 $H\subset G$ 为子群，$Y$ 为栖居于 $X$ 之上、带 $L^+H$-作用的空间。再设 $P_H$ 为 $X$ 上
主 $H$-丛。则 $Y$ 经 $P_H$ 的扭转是
\[
Y^{P_H}=Y\times^{L^+H}(P_H|_{D^\times})=(Y\times (P_H|_{D^\times}))/L^+H .
\]
本文把（上述构造的可因子化版本）施于 $\Gr_{G,\Ran}$，取 $H=T$，$P_T$ 为因子化
$T$-主丛
\[
\rho(\omega_X)=2\rho(\omega_X^{\otimes 1/2}).
\]
\srcpage{62}

由于 $LN$ 被 $L^+T$ 正规化，可取扭转 $LN_{\rho(\omega_X)}$。这是一个（因子化）群 ind-概形，
作用在 $\Gr_{G,\Ran,\rho(\omega_X)}$ 上。于是定义

\begin{definition}
\label{def:whit-G}
（带单位的）因子化范畴 $\Whit^{!}(G)\in \FactAlgCat$ 为
\[
\Whit^{!}(G)
=D\text{-}\mathrm{mod}_{1/2}\bigl(\Gr_{G,\rho(\omega_X)}\bigr)
^{(LN_{\rho(\omega_X)},\chi_{\rho(\omega_X)})}.
\]
\end{definition}

下面定义（$x_0$ 处的）因子化 $\Whit^{!}(G)$-模范畴 $\Whit^{!}(\Fl_G)$。先给出
$\Fl_{G,\Ran_{x_0}}$：对试验仿射概形 $S$ 及 $x\colon S\to \Ran_{x_0}$，$x$ 到
$\Fl_{G,\Ran_{x_0}}$ 的提升由四元组
\[
(P_G,\alpha,P_B,\gamma)
\]
给出：$S\times X$ 上的主 $G$-丛 $P_G$、在 $S\times X$ 中 $x$ 的图像之补集上 $P_G$ 的平凡化
$\alpha$，以及 $S\times\{x_0\}$ 上到 $B$ 的结构群约化
\[
\gamma\colon P_B\times^B G\xrightarrow{\ \sim\ }P_G|_{S\times\{x_0\}} .
\]
$\Fl_{G,\Ran_{x_0}}$ 自然带有 Ran Grassmann 叠上的因子化模空间结构。此外，因子化模群
ind-概形 $LG^{\Ran_{x_0}}$（即 $LG_{\Ran}$ 到 $\Ran_{x_0}\subset \Ran$ 的限制，
看作 $LG^{\Ran}$ 上因子化模空间中的群对象）作用在 $\Fl_{G,\Ran_{x_0}}$ 上。
于是可以仿照 $\Whit^{!}(G)$ 的定义作相应改动来定义 $\Whit^{!}(\Fl_G)$：作扭转
$\Fl_{G,\Ran_{x_0},\rho(\omega_X)}$，再取扭转 Whittaker 不变量
\[
D\text{-}\mathrm{mod}_{1/2}\bigl(\Fl_{G,\Ran_{x_0},\rho(\omega_X)}\bigr)
^{(LN_{\rho(\omega_X)},\chi_{\rho(\omega_X)})}.
\]
\srcpage{63}

\begin{definition}
\label{def:whit-Fl}
$\FactModCat_{x_0}$ 中的对 $(\Whit^{!}(G),\Whit^{!}(\Fl_G))$ 由因子化范畴
$\Whit^{!}(G)$ 与因子化 $\Whit^{!}(G)$-模范畴
\[
\Whit^{!}(\Fl_G)
=D\text{-}\mathrm{mod}_{1/2}\bigl(\Fl_{G,\Ran_{x_0},\rho(\omega_X)}\bigr)
^{(LN_{\rho(\omega_X)},\chi_{\rho(\omega_X)})}
\]
构成。
\end{definition}

下面转到谱侧。为定义因子化范畴 $\Rep(\check G)$ 与因子化 $\Rep(\check G)$-模范畴
$\QCoh(\check{\mathfrak{n}}/\check B)$，先回顾正则与亚纯局部系统的定义。
本小节余下内容的更详尽版本见 \cite[附录 B.3–B.7]{ABC+24a}。

对代数群 $H$，形式圆盘上正则 $H$-局部系统的因子化空间 $\LS^{\reg}_H$ 定义如下。

给定试验仿射概形 $S$ 及 $S$-点 $x\colon S\to \Ran$，$x$ 处的\emph{形式圆盘}定义为
$X$-概形
\[
\widehat{D}_x=(S\times X)\times_{S_{dR}\times X_{dR}}\Gamma_{x,dR};
\]
它下降到 $X_{dR}$-概形上，因为可以定义
\[
\widehat{D}_{x,\nabla}=(S\times X_{dR})\times_{S_{dR}\times X_{dR}}\Gamma_{x,dR}
=S\times_{S_{dR}}\Gamma_{x,dR},
\]
以及
\[
\widehat{D}_x=\widehat{D}_{x,\nabla}\times_{X_{dR}}X .
\]
设 $Y_\nabla$ 为 $X_{dR}$ 上的空间，记
\[
Y=Y_\nabla\times_{X_{dR}}X .
\]
\srcpage{64}

下面定义所谓的 $Y$ 的\emph{水平弧空间}，记作 $L^+_\nabla Y$。$x$ 到 $L^+_\nabla Y$ 的
提升由 $X_{dR}$-映射
\[
\widehat{D}_{x,\nabla}\to Y_\nabla
\]
给出。现在定义
\[
\LS^{\reg}_H=L^+_\nabla(\pt/H\times X).
\]
换言之，$x$ 到 $\LS^{\reg}_H$ 的提升由 $X_{dR}$-映射
\[
\widehat{D}_{x,\nabla}\to \pt/H\times X_{dR}
\]
给出，等价地由映射 $\widehat{D}_{x,\nabla}\to \pt/H$ 给出。

\begin{definition}
\label{def:rep}
因子化范畴 $\Rep(\check G)\in \FactAlgCat$ 定义为
\[
\QCoh(\LS^{\reg}_{\check G}).
\]
\end{definition}

形式圆盘上亚纯 $H$-局部系统的因子化空间 $\LS^{\mer}_H$ 定义起来更微妙。
对试验仿射概形 $S$ 及 $S$-点 $x\colon S\to \Ran$，注意 $\widehat{D}_x$ 是 ind-仿射
ind-概形。写成
\[
\widehat{D}_x=\colim \Spec(A_i),
\]
并记
\[
D_x=\Spec(\lim A_i),
\]
换言之，改在仿射概形的范畴中取余极限。于是 $x$ 处的\emph{去心形式圆盘}定义为
\[
D^\times_x=D_x\setminus \Gamma_x .
\]
\srcpage{65}

\begin{example}
\label{ex:formal-disk}
取 $S=\Spec(k)$，$x$ 为仅含 $x\in X$ 的单点集，$t$ 为 $X$ 在 $x$ 处的局部坐标。则
\[
\widehat{D}_x=\colim \Spec(k[t]/t^i)=\Spf(k[[t]]), \qquad
D_x=\Spec(\lim k[t]/t^i)=\Spec(k[[t]]).
\]
此时
\[
D^\times_x=\Spec(k((t))).
\]
\end{example}

由 \cite[附录 A.1]{Bog24}，$X$-概形 $D_x$ 与 $D^\times_x$ 都下降到 $X_{dR}$ 上。
把下降后的对象分别记作 $D_{x,\nabla}$ 与 $D^\times_{x,\nabla}$。若 $Y_\nabla\to X_{dR}$
仿射，则 $X_{dR}$-映射 $\widehat{D}_{x,\nabla}\to Y_\nabla$ 等价于 $X_{dR}$-映射
\[
D_{x,\nabla}\to Y_\nabla,
\]
于是可以把 $Y$ 的\emph{水平环路空间} $L_\nabla Y$ 定义为这样的因子化空间：$x$ 到
$L_\nabla Y$ 的提升由 $X_{dR}$-映射
\[
D^\times_{x,\nabla}\to Y_\nabla
\]
给出。
\srcpage{66}

\begin{remark}
\label{rem:loop-space-affine}
我们强调，$L_\nabla Y$ 只对仿射的 $Y_\nabla\to X_{dR}$ 定义。这是因为希望有一个由
$D^\times_{x,\nabla}\to D_{x,\nabla}$ 诱导的映射
\[
L^+_\nabla Y\to L_\nabla Y .
\]
正因如此，不能简单地把 $\LS^{\mer}_H$ 定义为 $L_\nabla(\pt/H\times X)$，
因为 $\pt/H\times X_{dR}$ 在 $X_{dR}$ 上并不仿射。
\end{remark}

下面引入 Jets 构造，见 \cite[附录 B.5]{ABC+24a}。对预叠 $Y\to X$，可作出 $X_{dR}$-预叠
\[
\Jets(Y)\to X_{dR},
\]
其 $S$-点由 $x\colon S\to X_{dR}$ 与 $X$-映射 $\widehat{D}_x\to Y$ 组成的对给出。
这就给出把 $Y$ 的因子化弧空间 $L^+Y$ 表示为
\[
L^+Y=L^+_\nabla\Jets(Y)
\]
的方式。此外，若 $Y\to X$ 仿射，则 $\Jets(Y)\to X_{dR}$ 也仿射，于是可以把 $Y$ 的
因子化环路空间 $LY$ 表示为
\[
LY=L_\nabla \Jets(Y).
\]
\srcpage{67}

现在可以定义 $\LS^{\mer}_H$ 了。为说明定义的动机，注意作为 $X_{dR}$ 上的群对象，
\[
H\times X_{dR}\subset \Jets(H\times X),
\]
而商
\[
\Jets(H\times X)/(H\times X_{dR})=\Jets(\mathfrak{h}\otimes \omega_X)
\]
（后者记作 $\Conn(\mathfrak{h})$）在 $X_{dR}$ 上仿射，因为 $\mathfrak{h}\otimes\omega_X$ 在
$X$ 上仿射（事实上是向量丛）。于是
\[
\pt/H\times X_{dR}=(\pt\times X_{dR})/(H\times X_{dR})
=\Conn(\mathfrak{h})/\Jets(H\times X),
\]
从而
\[
\LS^{\reg}_H=L^+_\nabla(\pt/H\times X)
=L^+_\nabla\Jets(\mathfrak{h}\otimes\omega_X)/L^+_\nabla\Jets(H\times X)
=L^+_\nabla\Conn(\mathfrak{h})/L^+H .
\]
定义
\[
\LS^{\mer}_H=L_\nabla\Conn(\mathfrak{h})/LH,
\]
即得所要的映射 $\LS^{\reg}_H\to \LS^{\mer}_H$。

本小节最后给出因子化 $\Rep(\check G)$-模范畴 $\QCoh(\check{\mathfrak{n}}/\check B)$
的定义。当然可以作因子化空间
\[
\LS^{\mer}_{\check B}\times_{\LS^{\mer}_{\check T}}\LS^{\reg}_{\check T},
\]
但我们要的是它在 $x_0$ 处、关于 $\LS^{\reg}_{\check G}$ 的因子化模空间版本。记
$\LS_H(D^{\times,\Ran}_{x_0})$ 为这样的因子化空间（定义见 \cite[定义 3.1.8]{Bog24}）：
$x\colon S\to \Ran_{x_0}$ 到 $\LS_H(D^{\times,\Ran}_{x_0})$ 的提升由
\[
\bigl\{D_{x,\nabla}\times_{X_{dR}}(X\setminus x_0)_{dR}\to \Conn(\mathfrak{h})\bigr\}
\big/\bigl\{D_{x,\nabla}\times_{X_{dR}}(X\setminus x_0)_{dR}\to H\bigr\}
\]
给出。由于 $D_{S\times x_0,\nabla}\subset D_{x,\nabla}$，经限制有映射
\[
\LS_H(D^{\times,\Ran}_{x_0})\to \LS_H(D^\times_{x_0}),
\]
后者是 $\LS^{\mer}_H$ 在 $\{x_0\}\in \Ran_{x_0}$ 处的纤维。
\srcpage{68}

把 $\check{\mathfrak{n}}/\check B^{\Ran}_{x_0}$ 定义为因子化 $\LS^{\reg}_{\check G}$-模空间
\[
\LS_{\check G}(D^{\times,\Ran}_{x_0})
\times_{\LS_{\check G}(D^\times_{x_0})}
\LS_{\check B}(D^\times_{x_0})
\times_{\LS_{\check T}(D^\times_{x_0})}
\LS_{\check T}(D_{x_0}).
\]

\begin{definition}
\label{def:rep-nB}
$\FactModCat_{x_0}$ 中的对 $(\Rep(\check G),\QCoh(\check{\mathfrak{n}}/\check B))$
由因子化范畴 $\Rep(\check G)$ 与因子化 $\Rep(\check G)$-模范畴
\[
\QCoh(\check{\mathfrak{n}}/\check B)
=\QCoh(\check{\mathfrak{n}}/\check B^{\Ran}_{x_0})
\]
构成。
\end{definition}
\srcpage{69}

\begin{remark}
\label{rem:nB-fiber}
$\check{\mathfrak{n}}/\check B^{\Ran}_{x_0}$ 在 $\{x_0\}\in \Ran_{x_0}$ 处的纤维是
\[
\LS_{\check B}(D^\times)\times_{\LS_{\check T}(D^\times)}\LS_{\check T}(D),
\]
它同构于光滑有限维叠
\[
\widetilde{\check N}/\check G=\check{\mathfrak{n}}/\check B;
\]
从而范畴层 $\mathrm{IndCoh}^{*}(\check{\mathfrak{n}}/\check B^{\Ran}_{x_0})$ 的纤维
正是通常的范畴 $\QCoh(\check{\mathfrak{n}}/\check B)$——这就是我们用同一记号
（滥用记号地）称呼该因子化模范畴的原因。
\yizhu{原文此处写作 $\mathrm{IndCoh}^{*}$，而定义~\ref{def:rep-nB} 中用的是
$\QCoh$，两处不一致。按上下文，此处所指应是 $\QCoh$，$\mathrm{IndCoh}^{*}$ 疑为笔误。}
{译者按（上下文推断，无外部依据）}
\end{remark}
